\SourceLeaf{231}%
\begingroup\brokenpenalty=10000\relax
\Article{6}
\indent Überträgt man die Richtung der Flächennormalen von der gekrümmten Fläche auf die Kugeloberfläche, so entspricht jedem bestimmten Punkt der gekrümmten Fläche ein bestimmter Punkt auf der Kugel. Ebenso wird jede Linie oder Figur auf der gekrümmten Fläche durch eine entsprechende Linie oder Figur auf der Kugel dargestellt. Beim Vergleich zweier Figuren, die sich auf diese Weise entsprechen, sodass die eine gleichsam ein Bild der anderen ist, sind zwei Gesichtspunkte zu beachten. Beim ersten betrachten wir allein die Größe, beim zweiten sehen wir von den Größenverhältnissen ab und betrachten allein die Lage.

Der erste Gesichtspunkt wird die Grundlage einiger Begriffe bilden, deren Aufnahme in die Lehre von den gekrümmten Flächen uns sinnvoll erscheint. Jedem von bestimmten Grenzen umschlossenen Teil einer gekrümmten Fläche schreiben wir eine \textit{Gesamt-} oder \textit{Integralkrümmung} zu, die durch den Flächeninhalt der entsprechenden Figur auf der Kugeloberfläche ausgedrückt wird. Von dieser Integralkrümmung ist sorgfältig eine Art spezifischer Krümmung zu unterscheiden, die wir das \textit{Krümmungsmaß} nennen werden. Es bezieht sich auf einen \textit{Punkt} der Fläche und bezeichnet den Quotienten, den man erhält, wenn man die Integralkrümmung des an den Punkt angrenzenden Flächenelements durch den Flächeninhalt des Elements selbst teilt. Es gibt damit das Verhältnis der einander entsprechenden unendlich kleinen Flächeninhalte auf der gekrümmten Fläche und auf der Kugeloberfläche an. Wir hoffen, dass sich der Nutzen dieser Neuerungen durch unsere späteren Ausführungen hinreichend erweisen wird. Was die Terminologie angeht, hielten wir es vor allem für wichtig, jede Mehrdeutigkeit auszuschließen. Daher hielten wir es für unangebracht, uns streng an die in der Lehre von den ebenen Kurven übliche, wenn auch nicht von allen gebilligte Terminologie anzulehnen. Nach dieser Analogie hätte man das Krümmungsmaß einfach Krümmung und die Integralkrümmung Amplitude nennen müssen. Warum aber sollten wir in der Wortwahl nicht frei sein, solange das Gemeinte Gehalt hat und die Bezeichnung keine falsche Auslegung nahelegt?

Die Lage einer Figur auf der Kugeloberfläche kann der Lage der entsprechenden Figur auf der gekrümmten Fläche ähnlich oder entgegengesetzt (umgekehrt) sein. Der erste Fall liegt vor, wenn zwei Linien, die auf der gekrümmten Fläche vom selben Punkt in verschiedenen, nicht entgegengesetzten Richtungen ausgehen, auf der Kugeloberfläche durch ähnlich liegende Linien dargestellt werden, etwa wenn das Bild der rechts liegenden Linie selbst rechts liegt. Der zweite Fall liegt vor, wenn das Gegenteil gilt. Wir werden diese beiden Fälle durch das positive oder negative \textit{Vorzeichen} des Krümmungsmaßes unterscheiden. Diese Unterscheidung ist aber offenbar nur möglich, wenn wir auf beiden Flächen %
